\chapter{Likelihood, profiling and calibrated inference}\label{ch:inference}

The marked likelihood provides the common bridge from structural and visibility questions to fitted inference. It records event times, pre-event intensities, the complete observation window, and the mark factors under the same mark model. Likelihood geometry can then be separated into three questions: whether a score direction has finite second moment, whether amplitude profiling gives a well-defined finite-history optimization, and how the rejection probability of a calibrated rule responds when the generating inhibition changes.

\section{The full marked likelihood and a score bound}

For an observed marked history \((T_i,Q_i)_{i=1}^{N(T)}\), write \(G_i=G_s(T_i-)\) and \(H_i=H_e(T_i-)\), so the event at \(T_i\) is excluded from both. With the common mark density \(p_Q\),
\begin{equation}
\ell(\vartheta,\eta)
=
\sum_{i=1}^{N(T)}
\log\left\{B(T_i)e^{-\eta G_i}+nH_i\right\}
-\int_0^T\left\{B(t)e^{-\eta G_s(t)}+nH_e(t)\right\}\dif t
+\sum_{i=1}^{N(T)}\log p_Q(Q_i).
\label{eq:full-likelihood}
\end{equation}
Every event term uses its pre-event stress and trigger history, the time integral covers the full horizon, and no first event is deleted. The mark sum belongs to the full likelihood. The first two terms form its event-time contribution, evaluated on the observed marked history; they are not the marginal time-only likelihood. The mark sum cancels in a contrast only when both candidates evaluate the same marks under the same mark law and support; it is constant in ground-rate coordinates when that mark law is held fixed. Ground-rate equality without equality of mark laws is insufficient for equality of marked likelihoods. The likelihood is conditioned on the declared starting history, rather than on the observed count.

For example, let a magnitude excess \(Q\ge0\) have exponential density \(\beta e^{-\beta Q}\), and let
\[
\kappa(Q)=(1-r)e^{r\beta Q},\qquad 0\le r<1.
\]
Then \(\E \kappa=1\), and the raw triggering amplitude is
\(n(1-r)e^{r\beta Q}\). Changing \(\beta\) while holding \(r\) fixed changes both the future mark distribution and the raw productivity exponent. To preserve a raw law \(C e^{cQ}\) on the observed history under a changed \(\beta'\), one must instead use
\[
r'=c/\beta',\qquad n'=C/(1-r'),
\]
subject to \(\beta'>c\) and the branching constraints. This can preserve the event-time likelihood contribution evaluated on the observed marked history; it does not preserve the full marked law. A mark-only estimate from a training prefix is also not automatically the joint marked-process maximum-likelihood estimate, since the ground-rate term can depend on \(\beta\) through productivity.

For fixed nuisance parameters, the inhibition score is
\begin{equation}
U_\eta=\ell_\eta
=\int_0^T
-\frac{B(t)e^{-\eta G_s(t)}G_s(t)}
{B(t)e^{-\eta G_s(t)}+nH_e(t)}
\left\{\dif N(t)-\lambda_\eta(t)\dif t\right\}.
\label{eq:eta-score}
\end{equation}
At \(\eta=0\), this is the score for the declared boundary direction, with the ordinary Hawkes law as reference. The next bound shows that raw heavy-tailed marks need not imply an infinite variance for this particular intensity-weighted score.

\begin{proposition}[Restricted finite score variance]
\label{prop:score_bound}
On \([0,T]\), with empty initial history and \(\eta\geq0\), let \(0<B(t)\le B_{\max}<\infty\), \(n>0\), let \(\kappa>0\) almost surely with \(\E \kappa=1\), and let \(h_e(u)>0\) for \(0\le u\le T\). Assume the same parents contribute to \(G_s\) and \(H_e\), with nonnegative \(g\) and \(\Phi_s\), finite mean event count, \(C_T=\int_0^T\Phi_s(u)^2/h_e(u)\,\dif u<\infty\), and \(m=\E [g(Q)^2/\kappa(Q)]<\infty\). Then the score \eqref{eq:eta-score} is square integrable and its predictable variance is
\[
\E _\eta U_\eta^2
=
\E _\eta\int_0^T
\frac{(B e^{-\eta G_s}G_s)^2}{\lambda_\eta}\,\dif t
\le
\frac{B_{\max}^2}{n}C_T\,\E _\eta N(T)m<\infty.
\]
\end{proposition}

\begin{proof}
When \(H_e(t)>0\), weighted Cauchy--Schwarz with weights \(\kappa(Q_i)h_e(t-T_i)\) gives
\[
G_s(t)^2
\le
H_e(t)\sum_{T_i<t}
\frac{g(Q_i)^2}{\kappa(Q_i)}
\frac{\Phi_s(t-T_i)^2}{h_e(t-T_i)}.
\]
When \(H_e=0\) there are no preceding parents and both displayed quantities are zero. Since \(\lambda_\eta\ge nH_e\) and \(B e^{-\eta G_s}\le B_{\max}\),
\[
\frac{(B e^{-\eta G_s}G_s)^2}{\lambda_\eta}
\le
\frac{B_{\max}^2}{n}
\sum_{T_i<t}
\frac{g(Q_i)^2}{\kappa(Q_i)}
\frac{\Phi_s(t-T_i)^2}{h_e(t-T_i)}.
\]
Integrating and applying Tonelli's theorem bounds the expected bracket by
\[
\frac{B_{\max}^2}{n}C_T
\E _\eta\sum_{T_i<T}\frac{g(Q_i)^2}{\kappa(Q_i)}.
\]
Fresh marks are drawn independently after acceptance, so the last expectation equals \(\E _\eta N(T)m\), although marks may influence future event counts. The stochastic-integral score is therefore an \(L^2\) martingale with the stated variance.
\end{proof}

Raw \(\E  g^2\) or \(\E \kappa^2\) may be infinite while the intensity-weighted ratio has finite moment. Positive \(n\) and matched triggering parents supply the required comparison history; the bound is not uniform as \(n\downarrow0\) or as the kernel-ratio integral grows.

\paragraph{Example at the synthetic mark parameters.}
For \(Q\sim\operatorname{Exp}(\beta)\),
\(g(Q)=e^{\gamma(Q-1)}\) and \(\kappa(Q)=(1-r)e^{r\beta Q}\),
where \(\gamma=\tfrac34\log 10\) is the local stress-mark exponent,
distinct from the auxiliary response gain in Chapter~\ref{ch:stability}.
The magnitude-six reference factor is included in \(g\); the outer
branching coefficient \(n\) is not included in \(\kappa\). Then
\[
  \E\!\left[\frac{g(Q)^2}{\kappa(Q)}\right]
  =\frac{e^{-2\gamma}}{1-r}
    \frac{\beta}{\beta(1+r)-2\gamma},
  \qquad \beta(1+r)>2\gamma.
\]
At the declared synthetic values \(\beta=2.849401764894274\) and \(r=0.7\),
the denominator is approximately \(1.390105\), giving
\(m\approx0.216065\). Although both first moments are finite,
\(2\gamma>\beta\) and \(2r\beta>\beta\), so
\(\E[g(Q)^2]=\E[\kappa(Q)^2]=\infty\).
The finite weighted moment verifies this mark condition in Proposition~\ref{prop:score_bound},
not all its other hypotheses by itself. The positive-triggering,
matched-parent, finite-horizon and kernel-ratio conditions remain essential to that bound;
the bound is not uniform near the simultaneous zero boundary, and
no joint Fisher-information, central-limit or Wilks conclusion follows.

At the simultaneous boundary \(n=\eta=0\), the conclusion can fail. Take a constant Poisson background \(B>0\), marks \(Q\sim\operatorname{Exp}(\beta)\), and \(g(Q)=e^{\gamma(Q-1)}\) with \(\gamma<\beta\le2\gamma\). The first stress-mark moment is finite while the second is infinite. For one exponential relaxation mode of rate \(k>0\), the resulting positive shot noise has a stationary version of mean \(B\E g/k\); copies with the same future points and marks have differences that decay exactly as \(e^{-kt}\). Thus a first-moment stable state process is compatible with the following score irregularity. On the positive-probability event of exactly one event at \(S<T\), the pre-event stress is zero and the score contribution has the form
\[
U_0=g(Q)\int_S^T B(t)\Phi_s(t-S)\,\dif t.
\]
The fresh mark is independent of that event configuration, and the integral is strictly positive on a positive-probability set of \(S\). Hence the score has infinite second moment. At positive \(\eta\), even with \(n=0\),
\[
\frac{(G_sB e^{-\eta G_s})^2}{\lambda_\eta}
\le B G_s^2e^{-\eta G_s}
\le\frac{4B}{e^2\eta^2},
\]
so integrable \(B\) controls the score bracket without a raw second mark moment. This regularization is not uniform as \(\eta\downarrow0\) and does not establish null-boundary asymptotic theory.

Finite score variance alone is not a central limit theorem. It also does not establish differentiability in quadratic mean for the full searched model or a regular information matrix at the boundary. The pathwise second derivative contains terms whose expectations may fail even when the martingale-score bracket is finite.

\section{Nuisance projection and local information}

Suppose the relevant scores are square integrable. Define the ensemble inner product
\[
\langle f,h\rangle_0
=
\E _0\int_0^T\lambda_0(t)f(t)h(t)\,\dif t.
\]
For a nuisance parameter \(\zeta\), the score integrand is
\(f_\zeta=\partial_\zeta\lambda_0/\lambda_0\). If the nuisance scores span a closed subspace, the residual after projecting \(f_\eta\) onto that subspace is the efficient-score integrand. In finite nonsingular notation,
\[
I_{\eta\mid\zeta}
=
I_{\eta\eta}-I_{\eta\zeta}I_{\zeta\zeta}^{-1}I_{\zeta\eta}
\le I_{\eta\eta}.
\]
When nuisance scores are redundant, a generalized inverse and its range condition replace the ordinary inverse.

This projection gives a quantitative meaning to masking by fitted background and triggering. Slowly varying depletion can resemble a background trend, while post-event depletion can resemble a triggering scale or productivity. It also shows why effective information is an ensemble object: it requires integrability and a reference law, not merely the algebra of one fixed path. At \(\eta=0\), the information contributed by an independently parameterized stress-kernel direction is zero, because
\(\partial_\theta\lambda=-\eta B e^{-\eta G_s}\partial_\theta G_s\).
An \(O(\eta^2)\) information statement away from zero requires a justified uniform moment bound for that derivative; the zero first derivative alone is not enough.

A related calculation can be made on one fixed reference path without the ensemble expectation. For \(B=\exp(a+dx)\), the four null log-intensity directions are
\[
s_a=\frac B{\lambda_0},\quad
s_d=\frac{xB}{\lambda_0},\quad
s_n=\frac{H_e}{\lambda_0},\quad
s_\eta=-\frac{BG_s}{\lambda_0}.
\]
The realized Gram matrix \(J_{ij}=\int_0^T\lambda_0s_is_j\,\dif t\) is positive semidefinite. Its residual inhibition norm after projecting on \((s_a,s_d,s_n)\) is the minimum of
\(\int\lambda_0(s_\eta-b_as_a-b_ds_d-b_ns_n)^2\dif t\).
If \(G_s=c_0+c_1x+c_2H_e/B\), that residual is exactly zero. This algebra explains how depletion can resemble a fitted level, trend and triggering component.

On two frozen reference paths, the archived residual/raw squared-norm ratios were approximately \(10.82\%\) and \(3.26\%\). These are two realized-path projections, rather than estimates of ensemble Fisher information. The paths and their calibrations differ; their ordering does not establish that greater exposure reduces expected information. Additional score directions can only reduce the residual for the same path and inner product, but no power decomposition follows from that projection fact.

\section{Conditional amplitude profiling}

Fix one observed marked history and fix the trend coefficient \(d\) and inhibition \(\eta\). Put
\[
q_i=\exp\{d x(T_i)-\eta G_i\},
\qquad
c_i=nH_i,
\qquad
J=\int_0^T\exp\{d x(t)-\eta G(t)\}\,\dif t,
\]
and let \(A=\exp(a)\) with \(A\in[L,U]\), where \(L=e^{-12}\) and \(U=1\). The negative event-time contribution to the joint marked log-likelihood, omitting the common mark constant and writing \(K\) for the full-window triggering compensator, is
\begin{equation}
F(A;d,\eta)=AJ+nK-\sum_{i=1}^{N}\log(Aq_i+c_i).
\label{eq:amplitude-objective}
\end{equation}
The full horizon and all pre-event intensities are retained. The mark factors in \eqref{eq:full-likelihood} are constant in these fitted coordinates under the common mark model.

\begin{proposition}[Conditional amplitude profile]
\label{prop:amplitude}
Assume a finite history and finite event features and parameters. Let every \(c_i\ge0\), let \(K\) be finite, and require \(0<J<\infty\). The specified finite-mode observation window satisfies these conditions when \(x\) is bounded and all recorded marks are finite. Then \eqref{eq:amplitude-objective} has a unique minimizer on \([L,U]\):
\begin{enumerate}
\item if \(F_A(L)\ge0\), the minimizer is \(L\);
\item otherwise, if \(F_A(U)\le0\), the minimizer is \(U\);
\item otherwise, the unique interior root of the continuous strictly increasing \(F_A\) is the minimizer.
\end{enumerate}
If \(N=0\), the unique minimizer is \(L\).
\end{proposition}

\begin{proof}
Ordinary finite-sum differentiation gives
\[
F_A=J-\sum_{i=1}^{N}\frac{q_i}{Aq_i+c_i},
\qquad
F_{AA}=\sum_{i=1}^{N}
\left(\frac{q_i}{Aq_i+c_i}\right)^2.
\]
For \(N\ge1\), \(q_i>0\), so \(F_{AA}>0\) and \(F_A\) is strictly increasing and continuous on the compact interval. If its endpoint values place zero inside the interval, it has one root, and the sign change from negative to positive identifies that root as the unique minimizer; endpoint cases place the minimizer at the corresponding boundary. If \(N=0\), then \(F(A)=AJ+nK\) is strictly increasing because \(J>0\), so the lower endpoint is the unique minimizer.
\end{proof}

The amplitude result is a finite-history optimization fact, not a global fit certificate. Convexity in \(A\) does not make the profiled surface \(V(d,\eta)=F(A_*(d,\eta);d,\eta)\) convex in \((d,\eta)\). A finite grid or local search establishes only its declared numerical search outcome.

For \(a=\log A\), set \(p_i=Aq_i/(Aq_i+c_i)\). Then
\[
g(a)=F_a=AJ-\sum_i p_i,
\qquad
g'(a)=AJ-\sum_i p_i(1-p_i),
\]
and \(g\) need not be monotone even though \(F_A\) is increasing in \(A\). A bracketing method for \(a\) can rely on continuity and endpoint signs, not on monotonicity of \(g\). At an interior root the exact balance is
\[
AJ=\sum_i p_i.
\]
The curvature identity in log amplitude is
\[
F_{aa}=\sum_i p_i^2+R,
\qquad
R=AJ-\sum_i p_i.
\]
Thus at the balanced root, the negative log-likelihood satisfies
\(F_{aa}=\sum_i p_i^2\), while the log-likelihood \(\ell=-F\) satisfies
\(\ell_{aa}=-\sum_i p_i^2\). At an approximate root, the balance residual \(R\) must be added to \(\sum_i p_i^2\); it is not exact curvature at that numerical point. In an interior fixed-\(d\) region, profile differentiation also gives
\[
V_d=AJ_x-\sum_i x(T_i)p_i,
\qquad
V_\eta=-AJ_G+\sum_i G_ip_i,
\]
where \(J_x\) and \(J_G\) integrate \(x(t)\) and \(G(t)\) against the background weight. Here \(J_x=\int_0^Tx(t)e^{dx(t)-\eta G(t)}\dif t\) and \(J_G=\int_0^TG(t)e^{dx(t)-\eta G(t)}\dif t\). At fixed \(d\), implicit differentiation in an interior-amplitude region gives

\[
\frac{\partial a_*}{\partial\eta}
=\frac{AJ_G-\sum_iG_ip_i(1-p_i)}{\sum_i p_i^2}.
\]
The numerator has no universal sign. An amplitude curve that also refits trend has a different total derivative. The identity \(AJ=\sum_i p_i\) equates integrated fitted background with event-wise background mixture weights; these weights are not observed physical parent labels or a stationary balance law. With empty prehistory and positive background, the first event has \(c_1=0\), so \(p_1=1\). No initial quiet interval or first event can be discarded from this argument.

These are local envelope and implicit formulas. Active endpoint transitions require one-sided treatment, and none of them implies that every starting point reaches the global profile minimum.

Stable evaluation also matters for the conditional proof. In logarithmic coordinates compute \(\log\lambda_i\) by a log-sum-exp of \(a+\log q_i\) and \(\log c_i\). When \(c_i=0\), the mathematically positive background gives \(p_i=1\), even if exponentiating \(\log q_i\) underflows. A numerical zero is not evidence that the exact objective is flat. The same amplitude algebra holds with positive quadrature weights, but agreement between quadrature orders at selected points is a numerical sensitivity check rather than a rigorous global integration bound.

\section{Boundary hypotheses and calibration}
\label{sec:boundary_calibration}

Feedback detection and testing equality of memories are different tasks. The first compares \(\eta=0\) with an active alternative; the second compares shared and separate stress/trigger families while permitting inhibition. At zero feedback, an independently varied stress-kernel parameter disappears from the observed law. Together with the one-sided constraint \(\eta\ge0\), this produces a nonregular null: there is a nuisance parameter present only under the alternative. A fixed-kernel boundary formula cannot automatically calibrate a maximization over stress kernels and fitted nuisances. The general difficulty is discussed by \citet{Davies1977}; the regularity conditions of boundary likelihood theory, including those in \citet{SelfLiang1987}, need their own verification here.

A calibration recipe must match the complete statistic actually computed. Under a fitted null, this means simulating new marked catalogues, refitting the allowed nuisance parameters under both candidates, repeating the declared kernel search and using the same failure and selection rules. A reference distribution for one nuisance block cannot be substituted for another. Such a plug-in procedure is a proposed finite-sample method; the present score bounds do not prove its uniform validity over a composite null or near weak identification. Non-rejection of shared memory cannot confirm that restriction without evidence that meaningful separate-memory alternatives are distinguishable.

The distinction between supplied and estimated nuisances is equally substantive. A likelihood ratio between two fully specified laws addresses a simple testing problem. A composite null can change its background and other permitted parameters, admitting alternative explanations that the simple null excludes. Superior calibrated detection with supplied generating values therefore does not measure purely avoidable fitting inefficiency. It combines different statistical questions with different procedures; neither an arbitrary optimizer improvement nor fixing real-data nuisance values without justification is guaranteed to recover the supplied-law performance.

\section{Response of a calibrated fixed-rule test}

Calibrated inference starts with a complete-data decision rule rather than with a particular fitted parameter value. Let \(I(\mathcal H)\in\{0,1\}\) be a fixed bounded measurable rejection indicator defined on every admissible marked history. The reference bank, tuning and model-selection procedure, critical threshold, treatment of empty histories, and any auxiliary randomization law are fixed independently of the true generating inhibition. The rule may estimate nuisances or select a model from the data, because these operations are part of the same measurable mapping from history to \(\{0,1\}\); what is held fixed is the rule as the generating \(\eta\) varies. Define
\[
p(\eta)=\E _\eta I(\mathcal H).
\]

\begin{proposition}[Positive-inhibition response of a fixed rule]
\label{prop:power_response}
Fix empty history, a common fresh mark law with finite nonnegative transforms, \(\E \kappa=1\), \(0\le n<1\), nonnegative kernels bounded on the finite window with normalized \(h_e\), and a deterministic nonnegative background with \(\mathcal B=\int_0^TB(t)\dif t<\infty\). Hold all these generating nuisances fixed as \(\eta\) varies. Then \(p(\eta)\) is continuously differentiable for \(\eta>0\), and
\[
p'(\eta)=\E _\eta[I(\mathcal H)S_\eta]
=\mathrm{Cov}_\eta(I(\mathcal H),S_\eta).
\]
Moreover,
\[
|p'(\eta)|
\le
\frac{2\sqrt{\mathcal B\,p(\eta)(1-p(\eta))}}{e\eta}
\le
\frac{\sqrt{\mathcal B}}{e\eta}.
\]
\end{proposition}

\begin{proof}
Let \(P_0\) be the noninhibited Hawkes law. Its mean count is at most \(\mathcal B/(1-n)\), and it dominates each inhibited construction. On a finite path the likelihood ratio is
\[
L_\eta=\exp\left\{\int_0^TB(1-e^{-\eta G_s})\dif t\right\}
\prod_i\frac{Be^{-\eta G_i}+nH_i}{B+nH_i}\le e^{\mathcal B}.
\]
At actual event times the reference intensity is positive almost surely; where both intensities vanish away from events, set their likelihood ratio to one. The stopped density martingale is uniformly bounded and hence defines the changed law. Its score \(S_\eta=\partial_\eta\log L_\eta=U_\eta\) satisfies
\[
\E _\eta S_\eta=0,\qquad
\E _\eta S_\eta^2\le\frac{4\mathcal B}{e^2\eta^2},
\]
by the positive-inhibition bracket bound, without a raw stress-mark second moment.

For clarity, write
\[
\begin{aligned}
 \mathcal B&=\int_0^T B(t)\,\dif t,
 &J_\eta&=\int_0^T B(t)(1-e^{-\eta G_s(t)})\,\dif t,\\
 q_i(\eta)&=\frac{B_i e^{-\eta G_i}+c_i}{B_i+c_i},
 &c_i&=nH_i\geq0,
\end{aligned}
\]
where $B_i=B(T_i)$ and $G_i,H_i$ are the pre-event features. Then
$L_\eta=e^{J_\eta}\prod_iq_i(\eta)$, with $0\leq q_i\leq1$ and
$0\leq J_\eta\leq\mathcal B$. At actual reference-law events
$B_i+c_i>0$ almost surely. On a compact positive interval
$\eta\in[\eta_-,\eta_+]$, direct differentiation gives
\[
 |q_i'(\eta)|
 =\frac{B_i}{B_i+c_i}G_i e^{-\eta G_i}
 \leq\frac{1}{e\eta_-},
 \qquad
 |J_\eta'|\leq\frac{\mathcal B}{e\eta_-}.
\]
The product rule, without division by any $q_i$, therefore yields
\[
\begin{aligned}
 |\partial_\eta L_\eta|
 &\leq e^{J_\eta}\left(
 |J_\eta'|\prod_iq_i+
 \sum_i |q_i'|\prod_{j\ne i}q_j\right)\\
 &\leq\frac{e^{\mathcal B}(\mathcal B+N(T))}{e\eta_-}.
\end{aligned}
\]
The same bound applies to $|\partial_\eta(L_\eta I)|$ for the fixed bounded
indicator. It is integrable under the reference law because
$\E _0N(T)\leq\mathcal B/(1-n)$. The bound is for the derivative of
$L_\eta$, not its logarithm. The finite-path derivative is continuous in $\eta$. Dominated convergence gives
\(p'(\eta)=\E _0[I\partial_\eta L_\eta]=\E _\eta[IS_\eta]\).
No derivative of the data-dependent fitted outputs is taken: \(I\) is the same mapping on every history. The mean-zero score gives the covariance identity. Cauchy--Schwarz with \(\mathrm{Var}(I)=p(1-p)\) gives the stated bounds.
\end{proof}

The covariance has no proved sign. Stronger generating inhibition need not increase detection probability for a particular calibrated rule. Positive \(\eta\) gives local differentiability and moment control, not monotonic power and not control as \(\eta\downarrow0\). Unknown, failed, capped, or otherwise unresolved computational outcomes remain unknown; they must not be recorded as nonrejections. Conditioning on completed cases estimates a different functional from the unconditional \(p(\eta)\) defined by the fixed rule.

The proposition also does not compare a simple hypothesis with a known generating law to a composite hypothesis whose nuisances are fitted. Those are different operating characteristics, with different permitted null explanations and decision rules. An oracle that knows the nuisance values and a composite rule that estimates them are not the same test function, so their performance gap is not by itself pure fitting waste.

\section{Scope and unresolved inference}

Taken together, the results separate four questions. Exact marked-law equality shows that zero feedback leaves a separately varied stress kernel unidentified; the restricted active-inhibition experiment can identify the inhibition parameter and kernel in an idealized full-law sense; the restricted score can remain square integrable even when raw mark moments diverge; and a calibrated fixed rule has a local covariance response at positive inhibition. Amplitude profiling adds a unique conditional optimization in \(A\), while leaving the remaining trend--inhibition surface without a global certificate.

These statements do not imply a Wilks law, a general Fisher-information theorem, likelihood-ratio asymptotics at \(\eta=0\), uniform validity near the boundary, a necessary detection threshold, or a finite-sample power result. They also provide no new numerical evidence and do not replace a specified calibration procedure by a general asymptotic claim. The unresolved inference is therefore explicit: structural identification in a restricted family is not practical detectability; finite score variance is not a distributional approximation; conditional amplitude uniqueness is not global identifiability; and a signed response formula for a fixed rule is not a theorem that stronger inhibition creates a more powerful test.
